%% ===========================================================================
\FloatBarrier
\section{Boundary conditions}
\label{sec:bc}
%% ===========================================================================

The fourth-order character of \eqref{eq:gov1} means that a boundary condition
on the displacement is not enough: the gradient theory requires a second,
higher-order condition on each edge, and the electrical condition decides
which of the two coupling factors the model returns. This section fixes both,
together with the axes that the rest of the paper refers to.

\FloatBarrier
\subsection{Mechanical and higher-order conditions}
\label{sec:bc_mech}

On an edge with unit outward normal $\bm{n}$ and tangent $\bm{t}$, stationarity
of \eqref{eq:energy} gives the classical traction
\begin{equation}
  t_i = \sigma_{ij}n_j - \partial_m\!\left(\varsigma_{ijm}\right)n_j
      + \partial_t\!\left(\varsigma_{ijm}\,t_m n_j\right),
  \label{eq:traction}
\end{equation}
and the higher-order double traction
\begin{equation}
  m_i = \varsigma_{ijm}\,n_j n_m ,
  \label{eq:doubletraction}
\end{equation}
which is conjugate to the normal derivative of the displacement. Three
combinations are used in the results:

\begin{enumerate}
\item \emph{clamped}: $u_i=0$ and $u_{i,j}n_j=0$ on the edge, so that both the
  displacement and its normal derivative -- the two quantities carried by the
  element \eqref{eq:dofs} -- are prescribed;
\item \emph{simply supported}: $u_i=0$ together with $m_i=0$, which leaves the
  slope free and removes the double force;
\item \emph{free}: $t_i=0$ and $m_i=0$, the natural conditions of the weak
  form, which are satisfied automatically by the discretization.
\end{enumerate}

The clamped condition is the one that produces the boundary layer of a
gradient theory: the gradient energy penalises the curvature that the
prescribed slope forces at the edge, so the solution relaxes over a length set
by the window. Section~\ref{sec:windowshape} measures that layer against
$a/L$, and the results are consistent with the ratio $a/L\le0.2$ that the
paper adopts as its validity window.

\FloatBarrier
\subsection{Electrical conditions: open and short circuit}
\label{sec:bc_elec}

Two electrical conditions are used. In the \emph{short-circuit} configuration
the electrode pair is grounded, the internal potential degrees of freedom are
condensed, and the mechanical problem has stiffness $K_{sc}$. In the
\emph{open-circuit} configuration one electrode is grounded and the other
floats at a single potential $V$ carrying zero net charge, which is again a
condensation and gives the rank-one update
\begin{equation}
  K_{oc} = K_{sc} + \frac{1}{\beta}\, d\,d^{\mathsf T},
  \label{eq:rankone}
\end{equation}
with $d$ the condensed coupling vector and $\beta$ the condensed dielectric
coefficient. Two coupling measures follow from this pair, and the difference
between them matters for the interpretation of every coupling number in the
paper:
\begin{enumerate}
\item the \emph{mode-based} factor of the $j$-th mode,
  $k_j^2=(\omega_{oc,j}^2-\omega_{sc,j}^2)/\omega_{oc,j}^2$, which depends on
  the eigenvector that the solver happened to return;
\item the \emph{mode-free quasi-static} factor of the electrode pair,
  \begin{equation}
    x = \frac{d^{\mathsf T}K_{uu}^{-1}d}{\beta},\qquad
    k^2 = \frac{x}{1+x},
    \label{eq:kstatic}
  \end{equation}
  which involves no eigenvector at all.
\end{enumerate}
Equation~\eqref{eq:kstatic} is the measure used wherever a quantity has to be
compared across geometries or across thicknesses, because it cannot select a
different mode in a different configuration. Four symbols carry the distinction
in what follows and are fixed here once: $k^2_{\rm qs}=x/(1+x)$ is the
mode-free factor of \eqref{eq:kstatic}; $k^2_{\rm qs,u}=x$ is the same
quantity without the coupled denominator, the one that can be set against
$e^{\mathsf T}C^{-1}e/\kappa$; $k^2_{j}$ is the mode-based factor of the
$j$-th mode; and $k^2_{j}(K_{uu})$ or $k^2_{j}(K_{sc})$ names the matrix whose
short-circuit eigenvalues enter it. An unsubscripted $k^2$ in the results is
$k^2_{\rm qs}$. The extensional coupling factor
of the \SI{200}{\micro\metre} by \SI{10}{\micro\metre} plate obtained from
\eqref{eq:kstatic} is $k^2=0.5581$, within $0.5\,\%$ of the constitutive
prediction $0.5608$ obtained from the constitutive tensors in
Sec.~\ref{sec:verification}.

\FloatBarrier
\subsection{Poling axes and Voigt indices}
\label{sec:axes}

The material is tetragonal $\mathrm{BaTiO_3}$ at room temperature, and the
crystal is anisotropic, so the two configurations of Fig.~\ref{fig:geometry}
do not share a modulus: poling along $x$ pairs the strain $\varepsilon_{xx}$
with $C^p_{xx}=9.5500\times10^{10}$\,Pa, while poling along $y$ pairs the same
strain with $C^p_{xx}=1.6946\times10^{11}$\,Pa, the two being exchanged by the
rotation that takes one poling axis into the other. Table~\ref{tab:axes} fixes
the axes, the Voigt indices and the tensors once, so that every later formula
can be read against it; a modulus may only be paired with the permittivity of
its own configuration, and the wrong pairing is examined, and excluded, in
Sec.~\ref{sec:verification}.

\begin{table}[!htbp]
\centering
\small
\caption{Which axis, which Voigt index, which modulus. (a) The axes of the
model: Voigt indices are $1=x$, $2=y$, $3=z$ throughout, and $z$ is a basal
direction of the $4mm$ crystal in both orientations. (b) The two
configurations used in the paper. They share the geometry and the material
record but not the poling axis, and a modulus may only be paired with the
permittivity of its own configuration.}
\label{tab:axes}
% table generated from the archived results
\begin{tabular}{@{}llll@{}}
\toprule
\multicolumn{4}{@{}l}{\emph{(a) the model axes, Voigt indices $1=x$, $2=y$, $3=z$}}\\
axis & role & extent & what it carries\\
\midrule
$x$ & length & \SI{200}{\micro\metre} & bending: $\varepsilon_{xx}=\kappa y$; the piezoelectric load direction\\
$y$ & thickness & \SI{10}{\micro\metre} & strain gradient $\partial\varepsilon_{xx}/\partial y=\kappa$; the flexoelectric gap\\
$z$ & out of plane & -- & neutral axis of the bending (curvature about $z$), free surface\\
\midrule
\multicolumn{4}{@{}l}{\emph{(b) the two configurations; they differ in the poling axis, and therefore in which entry of $\bm C_p$ multiplies $\varepsilon_{xx}$}}\\
& & piezo $k^2$ (\ref{sec:verification}) & flexo check (\ref{app:verification})\\
\midrule
poling axis & & $x$ & $y$\\
electrodes & & $x=0$, $x=L$ (short faces) & $y=0$, $y=h$ (long faces)\\
gap & & $L=$\SI{200}{\micro\metre} & $h=$\SI{10}{\micro\metre}\\
permittivity across the gap & & $\kappa_{xx}=\kappa_{33}=$\SI{4.9583e-10}{\farad\per\metre} & $\kappa_{yy}=\kappa_{33}=$\SI{4.9583e-10}{\farad\per\metre}\\
$C^p_{xx}$ & multiplies $\varepsilon_{xx}$ & \SI{9.5500e10}{\pascal} & \SI{1.6946e11}{\pascal}\\
$C^p_{yy}$ & & \SI{1.6946e11}{\pascal} & \SI{9.5500e10}{\pascal}\\
$E_b=1/(\bm C_p^{-1})_{xx}$ & beam modulus & \SI{7.6327e10}{\pascal} & \SI{1.3544e11}{\pascal}\\
coupling & & $\bm e_p^{\mathsf T}\bm C_p^{-1}\bm e_p/\kappa_{xx}$ & $3\mu^2/(C^p_{xx}h^2\kappa_{yy})$\\
\bottomrule
\end{tabular}

\end{table}

The pairing rule has a quantitative consequence that the paper relies on. The
extensional configuration is poled along $x$ with electrodes on the short
faces and the gap $L=\SI{200}{\micro\metre}$; the bending configuration is
poled along $y$ with electrodes on the long faces and the gap
$h=\SI{10}{\micro\metre}$. The two gaps differ by a factor of twenty, and so do
the electric fields produced by the same polarization, which is why the
extensional coupling is an electromechanical property of the material and the
bending coupling is a size effect of the geometry.

\FloatBarrier
